Chương này trình bày biểu đồ lớp (UML Class Diagram) cho các module cốt lõi của hệ thống IRMS, tuân thủ các nguyên lý SOLID.

\section{Module Quản lý Đơn hàng (Order Module)}
Đây là module trung tâm của hệ thống, quản lý vòng đời của một đơn hàng từ khi khởi tạo đến khi thanh toán xong. Các lớp chính bao gồm \texttt{Order}, \texttt{OrderItem} và \texttt{OrderService}.

\section{Module Quy trình Bếp (Kitchen Module)}
Module này chịu trách nhiệm điều phối việc chế biến món ăn. Lớp \texttt{KitchenService} phối hợp với \texttt{OrderPrioritizer} để sắp xếp thứ tự ưu tiên các món dựa trên thời gian chờ và loại món.

\section{Module Thanh toán (Payment Module)}
Quản lý các giao dịch tài chính. Lớp \texttt{PaymentProcessor} hỗ trợ nhiều chiến lược thanh toán khác nhau thông qua \textbf{Strategy Pattern}, cho phép dễ dàng tích hợp thêm các ví điện tử mới trong tương lai.

\section{Áp dụng nguyên tắc SOLID trong thiết kế}
\begin{itemize}
    \item \textbf{Single Responsibility (SRP)}: Mỗi service hoặc class chỉ đảm nhận một trách nhiệm duy nhất. Ví dụ: \texttt{InventoryService} chỉ tập trung vào quản lý kho, không tham gia vào logic thanh toán.
    \item \textbf{Open/Closed (OCP)}: Hệ thống mở cho việc mở rộng nhưng đóng cho việc sửa đổi. Việc sử dụng \textbf{Strategy Pattern} trong \texttt{PaymentModule} cho phép thêm phương thức thanh toán mới mà không cần sửa mã nguồn hiện có.
    \item \textbf{Liskov Substitution (LSP)}: Các lớp con hoặc implementation của Interface (như các loại \texttt{PaymentProvider} cụ thể) có thể thay thế hoàn toàn cho lớp cha/interface mà không làm thay đổi tính đúng đắn của hệ thống.
    \item \textbf{Interface Segregation (ISP)}: Thay vì một Interface lớn, hệ thống chia nhỏ thành các Interface chuyên biệt. Ví dụ: Một service chỉ cần xem dữ liệu đơn hàng sẽ dùng \texttt{OrderReader}, thay vì phải phụ thuộc vào một interface có cả chức năng chỉnh sửa hay xóa.
    \item \textbf{Dependency Inversion (DIP)}: Các module cấp cao không phụ thuộc vào module cấp thấp, cả hai đều phụ thuộc vào Abstraction. Các lớp nghiệp vụ phụ thuộc vào Interface thay vì các implementation cụ thể của tầng Database.
\end{itemize}

\section{Kết chương}
Các thiết kế chi tiết này là bản thiết kế kỹ thuật để đội ngũ lập trình có thể hiện thực hóa hệ thống một cách chính xác và đồng bộ. Chương tiếp theo sẽ tổng kết lại quá trình thực hiện và đưa ra hướng phát triển tương lai.
